% !TEX recipe = pdflatex
% Compilar: pdflatex -interaction=nonstopmode -halt-on-error topicos_micro.tex
% Executar duas vezes para atualizar o total de slides.
% Tema: https://latex-beamer.com/tutorials/beamer-themes/
% Conteúdo: ../../chapters/chap1/{0_introduction,1_literature_review}.tex
\documentclass[aspectratio=169,11pt]{beamer}
\usetheme{Madrid}
\usepackage[T1]{fontenc}
\usepackage[utf8]{inputenc}
\usepackage[brazilian]{babel}
\usepackage{lmodern}
\usepackage{booktabs,array,tabularx}
\definecolor{verde}{HTML}{24574D}
\definecolor{grafite}{HTML}{253A40}
\setbeamercolor{structure}{fg=verde}
\setbeamercolor{title}{bg=verde,fg=white}
\setbeamercolor{frametitle}{bg=verde,fg=white}
\setbeamercolor{normal text}{fg=grafite}
\setbeamertemplate{navigation symbols}{}
\setbeamertemplate{footline}{\leavevmode\hbox{\begin{beamercolorbox}[wd=.86\paperwidth,ht=2.7ex,dp=1.2ex,leftskip=1em]{author in head/foot}Capital natural e governança territorial\end{beamercolorbox}\begin{beamercolorbox}[wd=.14\paperwidth,ht=2.7ex,dp=1.2ex,center]{date in head/foot}\insertframenumber\,/\,\inserttotalframenumber\end{beamercolorbox}}}
\setbeamersize{text margin left=8mm,text margin right=8mm}
\setbeamertemplate{itemize items}[circle]
\newcommand{\fonte}[1]{\par\vfill{\fontsize{7.5}{9}\selectfont\color{gray}#1\par}}
\newcommand{\tpar}{\par\vspace{1.5mm}}
\newcolumntype{P}[1]{>{\raggedright\arraybackslash}p{#1}}
\title[Capital natural e governança territorial]{Capital Natural e Governança Territorial}
\subtitle{Produtividade, conversão e política ambiental em Mato Grosso}
\author{Fernando Souza da Vieira}
\institute{Tópicos de Microeconomia}
\date{}
\begin{document}

\begin{frame}
  \titlepage
\end{frame}

\section{Tema do trabalho}
\begin{frame}{Meio ambiente e decisões econômicas}
\begin{description}
\item[Meio ambiente] O estoque de vegetação integra as condições biofísicas que sustentam a produção agropecuária.
\medskip
\item[Endogenização] Decisões de conversão e conservação alteram esse estoque, que pode afetar a produtividade e os retornos econômicos.
\medskip
\item[Governança local] Regras, fiscalização e sanções condicionam os incentivos privados conforme as condições de cada território.
\medskip
\item[Intertemporalidade dos agentes] O uso da terra no presente altera estoques ambientais relevantes para oportunidades e riscos produtivos futuros.
\end{description}
\fonte{A intertemporalidade constitui um mecanismo conceitual do artigo, sem hipótese empírica independente. Fonte: introdução.}
\end{frame}

\section{Motivação}
\begin{frame}{Motivação: retornos privados e efeitos ambientais}
\begin{itemize}
\item A conversão pode ampliar a renda privada e reduzir serviços ecossistêmicos cujos benefícios não entram integralmente no cálculo do produtor.
\medskip
\item Alterações na vegetação afetam água, solo e microclima, com possíveis repercussões sobre a produtividade e o risco da atividade.
\medskip
\item A literatura oferece evidências para cada elo dessa relação, mas a identificação conjunta dos mecanismos permanece limitada.
\end{itemize}
\medskip
\begin{block}{Relevância econômica}
Compreender essa interação pode orientar decisões de uso da terra e a focalização territorial das políticas de conservação.
\end{block}
\fonte{Mann et al. (2012), Berman et al. (2023) e Redhead et al. (2020).}
\end{frame}

\begin{frame}{Mato Grosso no contexto da pesquisa}
\begin{itemize}
\item Produção agrícola e pecuária em larga escala, em um território que reúne \textbf{Amazônia, Cerrado e Pantanal}.
\medskip
\item Diferenças de vegetação, clima, disponibilidade hídrica e sistemas produtivos permitem investigar respostas heterogêneas.
\medskip
\item Custos de oportunidade e efetividade da fiscalização podem modificar a resposta dos agentes aos mesmos incentivos.
\medskip
\item A evidência sobre vegetação, choques climáticos e produtividade é desigual entre biomas, especialmente fragmentada para o Pantanal.
\end{itemize}
\fonte{VanWey et al. (2013), Simões et al. (2020), Sousa (2016), Seidl et al. (2001) e Cardoso et al. (2011).}
\end{frame}

\section{Problema}
\begin{frame}{Problema de pesquisa}
\begin{block}{Questão central}
Como retornos econômicos e instituições condicionam as decisões de conversão ou manutenção da vegetação, e em que medida o estoque de vegetação nativa altera a resposta da produtividade agrícola a choques climáticos em contextos ecologicamente e economicamente heterogêneos?
\end{block}
\medskip
\begin{description}
\item[Decisão de uso da terra] Examinar a interação entre retornos privados e condições institucionais.
\medskip
\item[Resposta produtiva] Investigar se a vegetação anterior ao choque modifica seus efeitos sobre a produtividade, com diferenças entre biomas.
\end{description}
\fonte{Pergunta de pesquisa reproduzida da introdução. As relações propostas constituem hipóteses a investigar.}
\end{frame}

\section{Hipóteses}
\begin{frame}{HG: interdependência entre economia e meio ambiente}
\begin{block}{Hipótese geral}
Retornos econômicos e instituições condicionam decisões de uso da terra que alteram o estoque ambiental, enquanto esse estoque pode, por sua vez, condicionar resultados produtivos relevantes para decisões econômicas subsequentes.
\end{block}
\medskip
\begin{itemize}
\item O mecanismo econômico-institucional explica a formação do estoque ambiental pela decisão de converter ou conservar.
\medskip
\item O mecanismo biofísico-produtivo relaciona esse estoque às condições de produção e aos retornos econômicos futuros.
\end{itemize}
\fonte{Formulação da introdução. H1 e H2 examinam os dois mecanismos que sustentam HG.}
\end{frame}

\begin{frame}{HG: respaldo na literatura}
\fontsize{8.5}{10.6}\selectfont
\setlength{\tabcolsep}{4pt}
\begin{tabularx}{\textwidth}{P{2.6cm} P{2.25cm} X P{2.25cm}}
\toprule
\textbf{Hipótese} & \textbf{Autores principais} & \textbf{Literatura} & \textbf{Outros autores} \\
\midrule
Retornos e instituições moldam o uso da terra e o estoque ambiental. Esse estoque pode afetar a produção e decisões econômicas posteriores.
& Angelsen (2010)\tpar Berman et al. (2023)\tpar Redhead et al. (2020)
& \textbf{Angelsen:} relaciona políticas de redução do desmatamento e produção agrícola. Fundamenta a conexão entre instituições, incentivos e conservação.\tpar
\textbf{Berman et al.:} relacionam preços agrícolas e desmatamento. Sustentam o elo entre retornos privados e alteração da cobertura vegetal.\tpar
\textbf{Redhead et al.:} relacionam composição e configuração da paisagem à resiliência dos rendimentos. Sustentam o elo entre condição ambiental e desempenho produtivo.
& Bradford et al. (2006)\tpar Upton et al. (2014)\tpar Mann et al. (2012)\tpar VanWey et al. (2013)\tpar Simões et al. (2020) \\
\bottomrule
\end{tabularx}
\fonte{As evidências respaldam elos do circuito H1 --> H2, mas não a identificação conjunta dos mecanismos. A heterogeneidade entre os três biomas permanece uma questão empírica do artigo.}
\end{frame}

\begin{frame}{H1: retornos, instituições e conversão}
\begin{block}{Hipótese específica 1}
A decisão de converter vegetação nativa para usos agropecuários responde à vantagem econômica esperada da conversão em relação à conservação, sendo essa resposta condicionada pelo ambiente institucional e podendo variar entre atividades, territórios e biomas.
\end{block}
\medskip
\begin{itemize}
\item \textbf{Vantagem econômica:} diferença entre os benefícios privados esperados da conversão e da conservação, considerados os custos relevantes.
\medskip
\item \textbf{Ambiente institucional:} direitos de propriedade, regras ambientais, monitoramento, fiscalização e sanções.
\end{itemize}
\fonte{Formulação e definições da introdução. O foco recai sobre a interação entre retornos e instituições e sua heterogeneidade territorial.}
\end{frame}

\begin{frame}{H1: respaldo na literatura}
\fontsize{8.5}{10.6}\selectfont
\setlength{\tabcolsep}{4pt}
\begin{tabularx}{\textwidth}{P{2.6cm} P{2.25cm} X P{2.25cm}}
\toprule
\textbf{Hipótese} & \textbf{Autores principais} & \textbf{Literatura} & \textbf{Outros autores} \\
\midrule
A conversão responde à vantagem econômica esperada, condicionada pelas instituições, com possíveis diferenças entre atividades, territórios e biomas.
& Harding et al. (2021)\tpar Assunção et al. (2023)\tpar Sousa (2016)
& \textbf{Harding et al.:} a lista prioritária reduziu o desmatamento associado aos preços das commodities. Respalda a interação entre pressão econômica e regulação.\tpar
\textbf{Assunção et al.:} identificam efeito do monitoramento por satélite e da fiscalização sobre o desmatamento. Sustentam o papel da efetividade institucional.\tpar
\textbf{Sousa:} encontra capacidade dissuasória dos embargos e baixa efetividade potencial das multas em Mato Grosso. Justifica distinguir instrumentos e condições locais.
& Angelsen (2010)\tpar Chomitz e Thomas (2003)\tpar Pacheco e Meyer (2022)\tpar Stickler et al. (2013)\tpar Bragança (2018) \\
\bottomrule
\end{tabularx}
\fonte{A heterogeneidade entre os três biomas permanece uma questão empírica do artigo.}
\end{frame}

\begin{frame}{H2: vegetação, clima e produtividade}
\begin{block}{Hipótese específica 2}
O estoque de vegetação previamente existente modifica, em magnitude economicamente relevante, a resposta da produtividade agrícola aos choques climáticos, podendo essa relação variar entre contextos produtivos, territórios e biomas.
\end{block}
\medskip
\begin{itemize}
\item \textbf{Anterioridade:} medir a cobertura vegetal antes do choque de precipitação, temperatura ou déficit hídrico.
\medskip
\item \textbf{Efeito de interesse:} diferenças na sensibilidade da produtividade ao choque conforme o estoque ambiental prévio.
\medskip
\item \textbf{Magnitude:} avaliar consequências substantivas para o desempenho produtivo, além da significância estatística.
\end{itemize}
\fonte{Formulação da introdução. H2 não pressupõe que maior vegetação implique maior produtividade média ou amortecimento em todos os contextos.}
\end{frame}

\begin{frame}{H2: respaldo na literatura}
\fontsize{8.5}{10.6}\selectfont
\setlength{\tabcolsep}{4pt}
\begin{tabularx}{\textwidth}{P{2.6cm} P{2.25cm} X P{2.25cm}}
\toprule
\textbf{Hipótese} & \textbf{Autores principais} & \textbf{Literatura} & \textbf{Outros autores} \\
\midrule
A vegetação anterior ao choque modifica a resposta da produtividade ao clima, em magnitude relevante, com possíveis diferenças entre territórios e biomas.
& Redhead et al. (2020)\tpar Spracklen et al. (2012)\tpar Dell et al. (2014)
& \textbf{Redhead et al.:} relacionam a composição e o arranjo da paisagem à resiliência dos rendimentos. Fundamentam a hipótese de sensibilidade produtiva diferenciada.\tpar
\textbf{Spracklen et al.:} associam a passagem do ar sobre florestas à precipitação tropical. Sustentam um mecanismo biofísico entre vegetação e condições de produção.\tpar
\textbf{Dell et al.:} discutem o uso de variações meteorológicas para estudar efeitos econômicos. Fundamentam a estratégia empírica, sem comprovar a moderação pela vegetação.
& Joshi e Garkoti (2023)\tpar Huang et al. (2021)\tpar Morel et al. (2019)\tpar Nóbrega et al. (2017)\tpar Xiao et al. (2023) \\
\bottomrule
\end{tabularx}
\fonte{Síntese da introdução e da revisão (RL11--RL14). A identificação do efeito da vegetação exige distingui-lo de solo, manejo e disponibilidade hídrica.}
\end{frame}
\end{document}

